\def\stat{zaharov}



\def\tit{К 40-ЛЕТИЮ ВЫХОДА ПОСТАНОВЛЕНИЙ ЦК КПСС И~СОВЕТА МИНИСТРОВ 

СССР О~ДАЛЬНЕЙШЕМ РАЗВИТИИ РАБОТ В~ОБЛАСТИ ВЫЧИСЛИТЕЛЬНОЙ 

ТЕХНИКИ, В~ТОМ~ЧИСЛЕ~В~АКАДЕМИИ НАУК СССР} %\\[-5pt]}



\def\titkol{К 40-летию выхода постановлений ЦК КПСС и~Совета Министров 

СССР} % о~дальнейшем развитии работ в~области вычислительной  техники, в~том числе в~АН СССР}



\def\aut{В.\,Н.~Захаров$^1$}



\def\autkol{В.\,Н.~Захаров}



\titel{\tit}{\aut}{\autkol}{\titkol}



\index{Захаров В.\,Н.}

\index{Zakharov V.\,N.}



%{\renewcommand{\thefootnote}{\fnsymbol{footnote}}

%\footnotetext[1]

%{Исследование выполнено с~использованием ЦКП <<Информатика>> ФИЦ ИУ РАН.}} 



\renewcommand{\thefootnote}{\arabic{footnote}}

\footnotetext[1]{Федеральный исследовательский центр <<Информатика и~управ\-ле\-ние>> Российской академии наук, 

\mbox{vzakharov@ipiran.ru}}



 \label{st\stat}



 \thispagestyle{headings}

 

 \vspace*{-14pt} %18pt

 

 



\Abst{В 2023~г.\ исполняется 40~лет с~момента образования Института проб\-лем 

информатики АН СССР, в~настоящее время входящего в~со\-став Федерального 

исследовательского цент\-ра <<Информатика и~управ\-ле\-ние>> РАН. Институт был 

образован в~соответствии с~По\-ста\-нов\-ле\-ни\-ем ЦК КПСС и~Совета Министров (СМ) \mbox{СССР} 

и~По\-ста\-нов\-ле\-ни\-ем Совета Министров  \mbox{СССР} от~29~июля 1983~г. В~статье 

рас\-смат\-ри\-ва\-ет\-ся некоторая предысто\-рия развития вычислительной техники в~СССР до 

выхода этих по\-ста\-нов\-лений, возвращение тематики работ по вы\-чис\-ли\-тель\-ной технике 

и~информатике в~АН \mbox{СССР}, создание Отделения информатики, вы\-чис\-ли\-тель\-ной техники 

и~автоматизации АН \mbox{СССР}, основные положения по\-ста\-нов\-ле\-ний. Приведены также краткие 

сведения о~развитии ИПИАН от момента создания до на\-сто\-яще\-го времени.}



%\vspace*{-8pt}

 

\KW{вычислительная техника; информатика; АН СССР; ЦК КПСС; Совет Министров 

\mbox{СССР}; Институт проб\-лем информатики}





%\vspace*{-8pt}



\DOI{10.14357\!/\!08696527230301}{XIWGAG}  



\vskip 10pt plus 9pt minus 6pt



 \thispagestyle{headings}

 

 \vspace*{-20pt}



\section{Введение}



\vspace*{-2pt}



   В СССР исследования в~области создания электронных вы\-чис\-ли\-тель\-ных 

машин (ЭВМ) начались в~конце 1940-х~гг.\ Пер\-вое в~\mbox{СССР} авторское свидетельство 

на изобретение автоматической циф\-ро\-вой вы\-чис\-ли\-тель\-ной машины (\mbox{АЦВМ}) 

на имя И.\,С.~Брука и~Б.\,И.~Рамеева было датировано 4~декабря 1948~г.\ 

(приоритет) и~выдано 16~февраля 1950~г. Создание пер\-вых в~\mbox{СССР} 

ЭВМ~--- М1 (И.~Брука) и~\mbox{МЭСМ} (малая электронная счетная машина)  (С.~Лебедева)~--- датируется 1951~г.~[1]. 

В~1950-е~гг.\ активно начались работы в~об\-ласти разработки новых моделей 

ЭВМ и~программного обеспечения для них. Пер\-вые отечественные машины 

создавались в~академических институтах. В~Москве работы велись под 

руководством чле\-на-кор\-рес\-пон\-ден\-та АН СССР И.\,С.~Брука сначала 

в~Энергетическом институте АН СССР (\mbox{ЭНИН}), где он организовал 

Лабораторию электросистем и~начал исследования в~об\-ласти расчетов 

мощ\-ных энергетических сис\-тем. В~1956~г.\ на базе Лаборатории  

элект\-ро\-сис\-тем \mbox{ЭНИН} АН \mbox{СССР} для разработки  

элект\-рон\-но-вы\-чис\-ли\-тель\-ной техники была образо-\linebreak\vspace*{-12pt}



\pagebreak



\noindent

вана 

самостоятельная Лаборатория управ\-ля\-ющих машин и~сис\-тем (ЛУМС) АН 

\mbox{СССР} под руководством И.\,С.~Бру\-ка. В~1957~г.\ он поставил 

научную проб\-ле\-му <<Разработка тео\-рии, принципов по\-стро\-ения 

и~применения специализированных вы\-чис\-ли\-тель\-ных и~управ\-ля\-ющих 

машин>>. Для ее решения~1~октября 1958~г.\ ЛУМС АН \mbox{СССР} была 

преобразована в~Институт электронных управ\-ля\-ющих машин (ИНЭУМ) 

АН СССР, директором которого стал И.\,С.~Брук.

   

   В Киеве под руководством академика АН УССР С.\,А.~Лебедева работы 

по созданию МЭСМ велись в~Лаборатории по спецмоделированию 

и~вы\-чис\-ли\-тель\-ной технике Института электротехники Академии наук 

Украинской ССР. В~1951~г.\ С.\,А.~Лебедев переехал в~Москву, где 

продолжил работу в~Институте точ\-ной механики и~вы\-чис\-ли\-тель\-ной техники 

АН СССР (ИТМиВТ), который воз\-глав\-лял с~1953 по 1973~гг.

   

   Кроме упомянутых двух институтов в~1950-е~гг.\ работы, связанные 

в~основ\-ном с~созданием средств программного обеспечения ЭВМ, велись 

в~целом ряде академических институтов, таких как Институт при\-клад\-ной 

математики, Вы\-чис\-ли\-тель\-ный центр АН СССР, Институт автоматики 

и~телемеханики АН СССР. 

   

   Работы по созданию аппаратных средств вы\-чис\-ли\-тель\-ной техники 

и~сис\-тем, базирующихся на этих средствах, проводились на предприятиях 

ряда промышленных министерств. Здесь мож\-но упомянуть Специальное 

конструкторское бюро №\,245 (СКБ-245), организованное в~1948~г., на базе 

которого в~1958~г.\ был образован НИИ электронных машин 

Минрадиопрома СССР. На\-уч\-но-про\-из\-вод\-ст\-вен\-ное предприятие 

(НПП) <<Рубин>> (г.~Пенза) было образовано По\-ста\-нов\-ле\-ни\-ем СМ 

\mbox{СССР} от 23.02.1953 как филиал Московского \mbox{СКБ-245}. 

По\-ста\-нов\-ле\-ни\-ем ЦК \mbox{КПСС} и~СМ \mbox{СССР} от 06.10.1958 филиал 

\mbox{СКБ-245} был преобразован в~НИИ управ\-ля\-ющих вы\-чис\-ли\-тель\-ных 

машин (\mbox{НИИУВМ}), а~через~8~лет Приказом Минрадиопрома 

\mbox{СССР} от 24.03.1966 переименован в~Пензенский НИИ 

математических машин. В~Москве работали завод счет\-ных аналитических 

машин (САМ) и~\mbox{НИИСчетмаш}. В~1957~г.\ Загорский электромеханический 

завод <<Звезда>> Минрадиопрома \mbox{СССР} был определен 

соразработчиком \mbox{ИТМиВТ}. В~1951~г.\ на правительственном уровне 

было принято решение о~строительстве Казанского завода математических 

машин (\mbox{КЗММ}), который вошел в~чис\-ло пер\-вых четырех заводов 

\mbox{СССР} по выпуску ЭВМ. С~1963~г.\ завод был подчинен 

Минрадиопрому \mbox{СССР} и~назван Казанским заводом электронных 

вычислительных машин (\mbox{КЗЭВМ}). В~1958~г.\ при Минском заводе 

математических машин Минрадиопрома \mbox{СССР} было создано 

Специальное конструкторское бюро (впоследствии~--- \mbox{НИИЭВМ}). 

В~Ленинграде в~1956~г.\ была создана Специальная лаборатория 

(\mbox{СЛ-11}) в~рамках \mbox{ОКБ-998} Авиапрома, позд\-нее 

превратившаяся в~Кон\-ст\-рук\-тор\-ско-тех\-но\-ло\-ги\-че\-ское бюро 

<<Свет\-ла\-на-мик\-ро\-элект\-ро\-ни\-ка>> (КБСМ). В~начале августа 

1954~г.\ был создан Вы\-чис\-ли\-тель\-ный центр №\,1 Министерства обороны, 

который спустя~7~лет был преобразован в~27~ЦНИИ МО \mbox{СССР}, 

сыграв\-ший значимую роль в~создании сис\-тем на базе средств 

вы\-чис\-ли\-тель\-ной техники, в~первую очередь в~военном применении.

   

   В начале 1960-х~гг.\ в~рамках так называемых <<хрущевских реформ>> 

руководство страны приняло решение, на\-прав\-лен\-ное на <<при\-бли\-же\-ние 

науки к~производству>>. Значительная часть институтов быв\-ше\-го Отделения 

технических наук АН СССР была передана в~про\-мыш\-лен\-ность, среди 

которых были ведущие институты в~об\-ласти вы\-чис\-ли\-тель\-ной техники 

(\mbox{ИНЭУМ}, \mbox{ИТМиВТ}). По\-сте\-пен\-но была изменена тематика 

этих институтов в~ущерб фундаментальным исследованиям 

и~перспективным разработкам. Изменилась и~мотивация ученых, 

работавших в~этих, теперь уже промышленных, институтах. Это решение 

привело к~значительному ослаб\-ле\-нию позиций АН СССР по 

на\-прав\-ле\-ни\-ям науки, связанным с~кибернетикой и~информатикой. Оно 

коснулось преж\-де всего Российской Федерации, академическая наука 

которой отож\-дест\-в\-ля\-лась с~АН \mbox{СССР}. Работы в~области кибернетики 

и~информатики стали развиваться в~это время в~Академиях наук союзных 

республик (Украины, Белоруссии, Эстонии, Латвии, Литвы, Армении, 

Грузии, Узбекистана), а~так\-же в~Сибирском и~Дальневосточном отделениях 

АН СССР, благодаря усилиям их организаторов~--- академиков 

М.\,А.~Лав\-ренть\-ева, С.\,Л.~Соболева, А.\,А.~Воронова~\cite{2-zah}. 

   

   В конце 1967~г.\ было принято Постановление ЦК \mbox{КПСС} и~Совета 

Министров \mbox{СССР} №\,1180-420 о~развитии средств производства 

вы\-чис\-ли\-тель\-ной техники в~стране, которым Минрадиопрому разрешалось 

создать На\-уч\-но-ис\-сле\-до\-ва\-тель\-ский центр электронной 

вы\-чис\-ли\-тель\-ной техники (\mbox{НИЦЭВТ}), который стал головным по 

созданию семейства ЭВМ <<Ряд>>, повторявших архитектуру машин IBM.



\section{Возвращение исследований в~области информатики 

и~вычислительной техники в~Академию наук}

   

   Как уже упоминалось выше, начало 1960-х~гг.\ было периодом, когда 

многие академические институты переводились в~промышленные 

министерства. Не избежал этого и~ИНЭУМ АН \mbox{СССР}~---  он был 

передан в~Госэкономсовет \mbox{СССР}, который в~дальнейшем стал 

Министерством при\-бо\-ро\-стро\-ения, средств автоматизации 

и~автоматизированных сис\-тем управ\-ле\-ния (Минприбор \mbox{СССР}).

{\looseness=1



}

   

   К началу 1980-х~гг.\ в~мире начался бум использования средств 

вычислительной техники буквально во всех областях. Это было вызвано 

в~значительной степени появлением на рынке нового вида массовой 

техники~--- персональных ЭВМ. В~\mbox{СССР} к~этому времени в~ряде 

министерств и~ведомств велись работы по разработке и~производству средств 

вы\-чис\-ли\-тель\-ной техники, однако уже отчетливо проявилось осознание 

заметного отставания от мировых лидеров в~этом на\-прав\-ле\-нии. В~конце 

1982~г.\ после смер\-ти руководителя страны Л.\,И.~Бреж\-не\-ва и~избрания на 

пост генерального секретаря ЦК КПСС Ю.\,В.~Андропова в~стране начались 

определенные изменения, за\-тро\-нув\-шие и~об\-ласть вы\-чис\-ли\-тель\-ной техники. 

Обновленным руководством страны был принят целый ряд важ\-ных решений, 

на\-прав\-лен\-ных на преодоление этого отставания. По правилам того времени 

такие решения оформлялись в~виде по\-ста\-нов\-ле\-ний ЦК \mbox{КПСС} или Совета 

Министров \mbox{СССР} или в~виде совместных по\-ста\-нов\-ле\-ний этих двух 

органов~\cite{3-zah}.

{\looseness=1



}



   

  В~1983~г.\  Общее собрание Академии наук \mbox{СССР} 

проходило~2~и~3~мар\-та. В своем вступительном слове Президент АН 

\mbox{СССР} А.\,П.~Александров говорил о~не\-об\-хо\-ди\-мости 

воз\-рож\-де\-ния и~усиления фундаментальных исследований в~об\-ласти 

компьютерных наук, а~так\-же це\-ле\-со\-образ\-ности (для обеспечения 

технологического паритета в~об\-ласти информационных технологий 

и~вы\-чис\-ли\-тель\-ной техники) создания специального отделения. С~большим 

докладом <<Об организации работ по информатике, вы\-чис\-ли\-тель\-ной технике 

и~автоматизации в~Академии наук \mbox{СССР}>> вы\-сту\-пил 

ви\-це-пре\-зи\-дент АН \mbox{СССР} Е.\,П.~Велихов. В~результате 

По\-ста\-нов\-ле\-ни\-ем Общего собрания АН \mbox{СССР} №\,12 от~3~марта 

1983~г.\ было принято решение об организации Отделения информатики, 

вы\-чис\-ли\-тель\-ной техники и~автоматизации (\mbox{ОИВТА}) в~со\-ста\-ве 

Секции фи\-зи\-ко-тех\-ни\-че\-ских и~математических наук при Президиуме 

АН \mbox{СССР}. Этим же по\-ста\-нов\-ле\-ни\-ем указанной секции было 

поручено представить на утверж\-де\-ние очередного Общего собрания АН 

\mbox{СССР} персональный со\-став Отделения.

{\looseness=1



}

   

   Распоряжением Президиума АН \mbox{СССР} от 24~ноября 1983~г.\ был 

утверж\-ден со\-став организационного бюро \mbox{ОИВТА} под 

председательством академика Е.\,П.~Велихова. В~него во\-шли: академики 

Ж.\,И.~Алферов, О.\,М.~Белоцерковский, В.\,А.~Мельников, 

чле\-ны-кор\-рес\-пон\-ден\-ты АН \mbox{СССР} К.\,А.~Валиев, А.\,П.~Ершов, 

Ч.\,В.~Копецкий, Б.\,Н.~Наумов, д.\,т.\,н.\ В.\,М.~Пономарев и~\mbox{к.\,ф.-м.\,н.} 

ученый секретарь Ю.\,С.~Виш\-ня\-ков.

   

   Постановлением Президиума АН СССР от 23~февраля 1984~г.\ были 

определены научные учреж\-де\-ния АН \mbox{СССР}, во\-шед\-шие в~со\-став 

\mbox{ОИВТА}. Это были четыре ранее существовавших института: 

Институт при\-клад\-ной математики им.\ М.\,В.~Келдыша, Вычислительный 

центр, Институт проб\-лем передачи информации и~Ленинградский  

на\-уч\-но-ис\-сле\-до\-ва\-тель\-ский вы\-чис\-ли\-тель\-ный центр, а~также 

пять новых институтов, созданных по\-ста\-нов\-ле\-ни\-ем ЦК КПСС и~СМ 

\mbox{СССР}, речь о~котором пойдет ниже. 

   

   Общее собрание АН СССР, состоявшееся 14~марта 1984~г., своим 

по\-ста\-нов\-ле\-ни\-ем №\,10 утвер\-ди\-ло персональный со\-став \mbox{ОИВТА}, 

в~который во\-шли 10~действительных членов АН \mbox{СССР} 

(Бе\-ло\-цер\-ков\-ский~О.\,М., Бун\-кин~Б.\,В., Ве\-ли\-хов~Е.\,П., Во\-ро\-нов~А.\,А., 

До\-род\-ни\-цын~А.\,А., Мель\-ни\-ков~В.\,А., Пугачев~В.\,С., Са\-мар\-ский~А.\,А., 

Семенихин~В.\,С.\ и~Тихонов~А.\,Н.) и~25~чле\-нов-кор\-рес\-пон\-ден\-тов 

(Алексеев~А.\,С., Би\-ца\-дзе~А.\,В., Бур\-цев~В.\,С., Валиев~К.\,А., 

Говорун~Н.\,Н., Гуляев~Ю.\,В., Ев\-ти\-хи\-ев~Н.\,Н., Емель\-я\-нов~С.\,В., 

Ер\-шов~А.\,П., Золотов~Е.\,В., Ко\-пец\-кий~Ч.\,В., Королев~Л.\,Н., 

Лавров~С.\,С., Лопато~Г.\,П., Макаров~И.\,М., Моисеев~Н.\,Н., 

Наумов~Б.\,Н., Попов~Е.\,П., Пос\-пе\-лов~Г.\,С., Ржа\-нов~А.\,В., Савин~А.\,И., 

Сифоров~В.\,И., Тихомиров~В.\,В., Цып\-кин~Я.\,З.\ и~Ше\-ре\-меть\-ев\-ский~Н.\,Н.).



%\vspace*{-9pt}



\section{Постановления ЦК КПСС и~Совета Министров СССР 

от~29~июля~1983~года}



\vspace*{-4pt}



   В этом году исполняется 40~лет выпущенным в~один день 29~июля 

1983~г.\ совместному Постановлению ЦК \mbox{КПСС} и~Совета Министров 

\mbox{СССР} №\,729-231 <<О~дальнейшем развитии работ в~об\-ласти 

вы\-чис\-ли\-тель\-ной техники>>~\cite{4-zah} и~детализирующему его 

По\-ста\-нов\-ле\-нию Совета Министров \mbox{СССР} №\,730-232 <<О~мерах по 

обеспечению работ в~об\-ласти вы\-чис\-ли\-тель\-ной техники и~ее применения 

в~народном хозяйстве>>~[5], сыграв\-ших большую роль в~истории 

информатики в~стране (см.\ рисунок).

   

    \begin{figure}[b]

    \vspace*{-6pt}

  \begin{center}

 \mbox{%

 \epsfxsize=126mm 

\epsfbox{zah-1.eps}

 }

\vspace*{4pt}



    {\small Копии первых страниц постановлений ЦК \mbox{КПСС} и~СМ \mbox{СССР}}

    \end{center}

    \vspace*{-4pt}

    \end{figure}

   

   В постановлении ЦК КПСС и~СМ \mbox{СССР}, под которым стоят 

подписи секретаря ЦК КПСС Ю.~Анд\-ро\-по\-ва и~Председателя СМ 

\mbox{СССР} Н.~Тихонова, были сформулированы глав\-ные задачи  

в~об\-ласти исследований, разработки, производства и~применения 

вы\-чис\-ли\-тель\-ной техники в~народном хозяйстве на ближайшую перспективу:\\[-15pt]

   \begin{itemize}

   \item проведение фундаментальных и~при\-клад\-ных исследований по 

созданию средств вы\-чис\-ли\-тель\-ной техники на основе перспективной  

кон\-ст\-рук\-тив\-но-эле\-мент\-ной базы, создание новых 

высокопроизводительных вы\-чис\-ли\-тель\-ных и~управ\-ля\-ющих комплексов 

(в~том числе специализированных) и~сетей цент\-ров общего пользования, 

расширение областей эффективного применения средств вы\-чис\-ли\-тель\-ной 

техники;\\[-13.5pt]

   \item разработка и~освоение производства технических, а~так\-же 

программных средств вы\-чис\-ли\-тель\-ной техники и~носителей информации, 

со\-от\-вет\-ст\-ву\-ющих по техническому уров\-ню лучшим мировым образцам;\\[-13.5pt]

   \item повышение на\-деж\-ности и~совершенствование структуры выпуска 

средств вы\-чис\-ли\-тель\-ной техники, унификация и~стандартизация 

технических, конструктивных и~программных решений, улуч\-ше\-ние 

обеспечения вы\-чис\-ли\-тель\-ных и~управ\-ля\-ющих комплексов 

периферийными устройствами и~по\-став\-ка этих комплексов в~конфигурациях, 

со\-от\-вет\-ст\-ву\-ющих требованиям пользователей вы\-чис\-ли\-тель\-ной техники;\\[-13.5pt]

   \item дальнейшее развитие и~интеграция автоматизированных сис\-тем 

различного уров\-ня и~назначения, а~так\-же повышение экономической  

эф\-фек\-тив\-ности их в~народном хозяйстве;\\[-13.5pt]

   \item создание сис\-тем передачи и~телеобработки данных для сетей 

вы\-чис\-ли\-тель\-ных цент\-ров коллективного пользования (ВЦКП);\\[-13.5pt]

   \item широкое применение микропроцессоров и~мик\-ро-ЭВМ при 

одновременном снижении их сто\-и\-мости;\\[-13.5pt]

   \item организация разработки и~производства на промышленной основе 

про\-грам\-мных средств для вы\-чис\-ли\-тель\-ных и~управ\-ля\-ющих комплексов 

и~автоматизированных сис\-тем;\\[-13.5pt]

   \item подготовка специалистов по разработке, обслуживанию 

и~применению технических и~программных средств вы\-чис\-ли\-тель\-ной 

техники;\\[-13.5pt]

   \item расширение круга пользователей вы\-чис\-ли\-тель\-ной техникой путем 

организации сис\-те\-мы переподготовки специалистов во всех отраслях 

народного хозяйства, включая руководящих работников.

   \end{itemize}

   

   На Академию наук СССР этим постановлением была возложена 

ответственность за проведение фундаментальных и~прикладных 

исследований в~об\-ласти вы\-чис\-ли\-тель\-ной техники и~автоматизированных 

систем, а~так\-же научное руководство такими работами, проводимыми 

организациями Академий наук союзных рес\-пуб\-лик, выс\-шей школы, 

министерств и~ведомств. Было предписано считать это одной из важ\-ней\-ших 

задач Академии наук \mbox{СССР}. 

   

   Именно в~этом совместном постановлении ЦК \mbox{КПСС} и~СМ \mbox{СССР} 

было сказано: <<Принять предложение АН СССР, согласованное с~ГКНТ 

и~комиссией Президиума СМ \mbox{СССР} по во\-ен\-но-про\-мыш\-лен\-ным вопросам 

о~создании в~сис\-те\-ме АН \mbox{СССР}: Научного центра 

по фундаментальным проб\-ле\-мам вы\-чис\-ли\-тель\-ной техники и~сис\-тем  

управ\-ле\-ния (с~включением в~его со\-став организуемых в~г.~Яро\-слав\-ле 

Института проб\-лем вы\-чис\-ли\-тел\-ьной техники, Института 

микроэлектроники, СКБ и~опыт\-но\-го производства) на правах  

на\-уч\-но-тех\-ни\-че\-ского объединения; Института проб\-лем кибернетики 

на базе лабораторий научного совета АН СССР по комплексной проб\-ле\-ме 

<<Кибернетика>> в~г.~Моск\-ве с~филиалом  

в~г.~Пе\-ре\-славль-За\-лес\-ский Ярославской об\-ласти; Института  

проб\-лем технологии мик\-ро\-элект\-ро\-ни\-ки и~особо чис\-тых материалов (со 

специальным конст\-рук\-торско-тех\-но\-ло\-ги\-че\-ским бюро и~опыт\-ным 

производством) в~пос.~Чер\-но\-го\-лов\-ка на базе ряда подразделений 

Института физики твер\-до\-го тела и~других организаций АН \mbox{СССР}. 

Разрешить АН \mbox{СССР} создать Институт проб\-лем информатики 

с~опыт\-ным производством в~г.~Моск\-ве и~с~филиалами в~гг.~Казани 

и~Бердянске За\-по\-рож\-ской об\-ласти>>. Эти созданные организации 

в~1984~г.\ во\-шли в~со\-став \mbox{ОИВТА}.

    

   В совместном постановлении были утверж\-де\-ны <<Основ\-ные 

на\-прав\-ле\-ния фундаментальных и~при\-клад\-ных исследований в~об\-ласти 

вы\-чис\-ли\-тель\-ной техники и~автоматизированных сис\-тем на период до 

1990~года>>. Приведем их здесь пол\-ностью для того, чтобы почувствовать 

<<аромат эпохи>>.



\vspace*{-3pt}



\subsection*{ЭВМ и~вычислительные системы}



\vspace*{-3pt}



     Научная разработка единой стратегии развития вы\-чис\-ли\-тель\-ных 

и~автоматизированных управ\-ля\-ющих сис\-тем на основе унификации 

и~стандартизации технических и~программных средств.

     

     Создание ЭВМ сверхвысокой производительности (свыше~1~млрд 

операций в~секунду) с~использованием новых принципов организации 

вы\-чис\-ли\-тель\-но\-го процесса и~по\-стро\-ения архитектуры ЭВМ.

     

     Разработка новых принципов создания специализированных 

мик\-ро\-про\-цес\-со\-ров, микро- и~ми\-ни-ЭВМ.

     

     Разработка перспективных оперативных и~внеш\-них  

за\-по\-ми\-на\-ющих устройств (ЗУ) и~широкого набора периферийных средств.

     

     Исследование и~разработка перспективных технических средств 

вычислительной техники с~использованием новых физических принципов 

(криоэлектроники, оптоэлектроники, акусто- и~маг\-ни\-то\-элект\-ро\-ни\-ки).

     

     Разработка языков формализованного описания архитектуры ЭВМ для 

сис\-тем автоматизированного проектирования машин.

     

     Исследования в~об\-ласти искусственного интеллекта.

     

     \vspace*{-3pt}

     

     \subsection*{Программное обеспечение}

     

     \vspace*{-3pt}

     

     Исследование и~создание методов математического моделирования  

раз\-ра\-ба\-ты\-ва\-емых ЭВМ.

     

     Создание современной технологии разработки программного 

обеспечения, включая разработку методов автоматизированного  

по\-стро\-ения программных сис\-тем.

     

     Исследование и~разработка машинонезависимых операционных  

сис\-тем.

     

     Исследование и~разработка языков программирования высокого  

уров\-ня, в~том чис\-ле логических (непроцедурных), метаязыков.



\vspace*{-6pt}

     

     \subsection*{Микроэлектронная база}

     

     \vspace*{-3pt}

     

     Исследование физических и~технологических проб\-лем создания 

суб\-мик\-рон\-ных структур на перспективных материалах, создание опытных 

образцов технологического оборудования.

     

     Разработка новых принципов создания сверх\-быст\-ро\-дей\-ст\-ву\-ющих 

и~сверхбольших интегральных схем и~их макетирование.

     

     Разработка новых методов контроля полупроводниковых материалов 

и~структур сверх\-быст\-ро\-дей\-ст\-ву\-ющих и~сверхбольших интегральных схем.

     

     Разработка технологии изготовления специальных и~особо чис\-тых 

материалов и~совершенных крис\-тал\-лов.

     

     Разработка методов и~средств контроля технологических процессов 

и~качества изделий мик\-ро\-элект\-рон\-ной базы.

     

     Изучение факторов, опре\-де\-ля\-ющих долговременную ста\-биль\-ность 

элементной базы мик\-ро\-элект\-ро\-ники.

     

     Создание банков данных для разработки и~оптимизации процессов 

из\-го\-тов\-ле\-ния особо чис\-тых материалов, сверх\-быст\-ро\-дей\-ст\-ву\-ющих 

и~сверхбольших интегральных схем.



\vspace*{-6pt}

     

     \subsection*{Применение вычислительной техники}

     

     \vspace*{-3pt}

     

     Исследование принципов и~основных на\-прав\-ле\-ний создания сетей 

ЭВМ и~\mbox{ВЦКП}.

     

     Разработка методологии создания комплексных интегрированных 

автоматизированных сис\-тем, охватывающей все этапы разработки 

и~из\-го\-тов\-ле\-ния этих сис\-тем.

     

     Разработка теоретических основ и~методологии по\-стро\-ения 

взаимоувязанных автоматизированных сис\-тем различных уров\-ней  

и~на\-прав\-ле\-ний.

   

   В Постановлении СМ СССР~[5], подписанном Председателем СМ 

\mbox{СССР} Н.~Тихоновым и~управ\-ля\-ющим делами СМ \mbox{СССР} 

М.~Смир\-тю\-ко\-вым, были определены конкретные задачи по проведению 

фундаментальных и~при\-клад\-ных исследований, пе\-ре\-чис\-ле\-ны задания на 

разработку технических заданий, изготовление опыт\-ных образцов и~освоение 

серийного производства новых средств вы\-чис\-ли\-тель\-ной техники на 

несколько ближайших лет. Этим же по\-ста\-нов\-ле\-ни\-ем предусмат\-ри\-вал\-ся 

целый комплекс мер, необходимых для реализации по\-став\-лен\-ных задач. Так, 

было предписано осуществить строительство, расширение и~реконструкцию 

предприятий и~организаций по спис\-ку с~выделением средств, разработку 

программы стандартизации и~межотраслевой унификации средств 

вы\-чис\-ли\-тель\-ной техники, создание хозрасчетных вы\-чис\-ли\-тель\-ных цент\-ров. 

Были и~такие непривычно сейчас звучащие слова: <<ГКНТ, Комиссии по 

ВПК и~АН \mbox{СССР} осуществить перевод в~6-ме\-сяч\-ный срок 

до~300~квалифицированных специалистов из МРП, Минприбора и~МЭП 

в~ИПИ АН. Разрешить АН \mbox{СССР} раз\-мес\-тить ИПИ АН в~здании 

Президиума АН \mbox{СССР}, строительство которого предусмот\-ре\-но 

распоряжением СМ \mbox{СССР} от 01.12.1975 №\,2614. Увеличить на 

1983~г.\ АН \mbox{СССР} лимит чис\-лен\-ности работников организаций, 

расположенных в~г.~Москве, на 100~человек>>.

   

   С исторической точки зрения представляют интерес задания по созданию 

средств вы\-чис\-ли\-тель\-ной техники, содержащие конкретные па\-ра\-мет\-ры, 

ха\-рак\-те\-ри\-зу\-ющие уровень того времени и~ближайшие перспективы. 

Приведем лишь некоторые примеры. 

   

   Так, ставилась задача по исследованию принципов создания накопителей 

на магнитных дис\-ках (НМД) ем\-костью до~3000~МБайт, со ско\-ростью 

обмена данными до~3000~КБайт\!/\!с, сред\-ним временем до\-сту\-па~20--25~мс; 

по разработке методов создания ЗУ на цилиндрических магнитных доменах 

(ЦМД) ем\-костью~10$^8$~бит на крис\-талл; по исследованию принципов 

создания оп\-ти\-ко-ме\-ха\-ни\-че\-ско\-го дискового ЗУ ем\-костью 

$(3\mbox{--}5)\cdot 10^{12}$~бит, со ско\-ростью обмена данными 

до~2~МБайт\!/\!с, средним временем доступа~15~с. Ставилась задача по 

производству в~1985~г.\ накопителей на гибких магнитных дисках

 ем\-костью~0,14~МБайт, ско\-ростью 

обмена~250~кбит\!/\!с, с~плот\-ностью записи~109\!/\!218~бит\!/\!мм, диа\-мет\-ром 

дис\-ка~130~мм (НИИ периферийного оборудования, Киев). В~час\-ти 

мейнфреймов в~задании на~1984~г.\ фигурировала ЭВМ \mbox{ЕС-1046} 

 про\-из\-во\-ди\-тель\-ностью~1~млн~оп\!/\!с, ем\-костью ОЗУ до~8~МБайт 

(разработка и~изготовление опыт\-но\-го образца \mbox{НИЦЭВТ} 

и~\mbox{ЕРНИИММ}, серийное производство~--- Казанский завод ЭВМ 

МРП). А~на 1988~г.\ было задано уже производство ЭВМ \mbox{ЕС-1087}  

про\-из\-во\-ди\-тель\-ностью~18~млн~оп\!/\!с, ем\-костью ОЗУ до~128~МБайт 

(Минское производственное объединение вычислительной техники, МПО ВТ). В~час\-ти мини- и~мик\-ро\-ЭВМ говорилось о~за\-пус\-ке 

в~производство в~1984~г.\ вы\-чис\-ли\-тель\-но\-го комплекса \mbox{СМ-1420}, 

программно совместимого с~ЭВМ \mbox{СМ-3} и~\mbox{СМ-4}, 

раз\-ряд\-ностью~16~двоичных разрядов, про\-из\-во\-ди\-тель\-ностью до~1~млн~оп\!/\!с,  

ем\-костью ОЗУ~2~МБайт (ИНЭУМ и~Киевское ПО <<Электронмаш>> им.\ 

В.\,И.~Ленина Минприбора), и~микроЭВМ <<Элект\-ро\-ни\-ка-60-1>> 

раз\-ряд\-ностью~16~двоичных разрядов, про\-из\-во\-ди\-тель\-ностью до~600~тыс.~оп\!/\!с, ем\-костью ОЗУ до~256~КБайт (ПО <<Электроника>> 

Мин\-электрон\-про\-ма).

   

   О масштабах проводимых работ можно судить, например, по перечню, 

содержащемуся в~задании по строительству новых, расширению 

и~реконструкции дей\-ст\-ву\-ющих предприятий и~организаций, 

задействованных в~выполнении по\-ста\-нов\-ле\-ния. Приведем его  

пол\-ностью.

   

   \textbf{Минрадиопром:} НИЦЭВТ (Москва), Загорский 

электромеханический завод (Московская об\-ласть), Костромской 

электромеханический завод, завод САМ (Москва), Волжский завод 

радиотехнических элементов (Волгоградская об\-ласть), Астраханский 

машиностроительный завод <<Прогресс>>, Пензенский завод <<ВЭМ>>, 

Минский завод ЭВМ, Минский завод многослойных печатных плат, 

Кишиневский завод счетных машин, Тбилисский завод  

элект\-рон\-но-вы\-чис\-ли\-тель\-ной техники, Ташкентский завод 

электронных вы\-чис\-ли\-тель\-ных машин <<Алгоритм>>, Кировский 

приборостроительный завод, НИИ ЭВМ (Минск), Казанский завод пишущих 

устройств, Казанский завод ЭВМ, Боярский машиностроительный завод 

<<Искра>> (Киевская об\-ласть), Брестский электромеханический завод, 

завод <<Электроприбор>>  

(г.~Ка\-ме\-нец-По\-доль\-ский Хмельницкой об\-ласти), Каневский 

электромеханический завод <<Магнит>> (Черкасская об\-ласть), 

Фрунзенский завод ЭВМ, НИИ вы\-чис\-ли\-тель\-ной техники (Пенза), Бакинский 

завод ЭВМ, Бакинский радиозавод, Казанский НИТИ вы\-чис\-ли\-тель\-ной 

техники, КБ <<Север>> (Киров), ПО <<Центр\-ЭВМ\-комп\-лекс>> (Москва).

   

   \textbf{Минприбор:} Орловский завод УВМ, Киевский завод электронных 

вычислительных и~управ\-ля\-ющих машин, Северодонецкий 

приборостроительный завод, Вильнюсский завод счетных машин, Рязанский 

завод САМ, Таурагский завод элементов вы\-чис\-ли\-тель\-ных машин (Литовская 

ССР), Паневежский завод точ\-ной механики (Литовская ССР), Курский завод 

<<Счетмаш>>.

   

   \textbf{Минэлектронпром:} завод <<Альтаир>> (Ярославль), завод 

<<Процессор>> (Воронеж), ПО <<Электроника>> (Воронеж), Брянский 

завод полупроводниковых приборов, Кишиневский завод <<Мезон>>, НИИ 

физических проб\-лем (Зеленоград), НИИ молекулярной электроники 

(Зеленоград), ЦНИИ <<Цик\-лон>> (Москва), завод микропроцессоров 

(Бельцы Кишиневской об\-ласти), завод <<Девиз>> (г.~Алексеевка 

Белгородской об\-ласти), Калининградский машиностроительный завод, 

Выборгский приборостроительный завод, Черняховский 

машиностроительный завод (Калининградская об\-ласть).

   

   \textbf{АН СССР:} Научный центр АН СССР (Ярославль), филиал 

Института кибернетики АН \mbox{СССР} (Пе\-ре\-славль-За\-лес\-ский), Институт 

проб\-лем кибернетики АН \mbox{СССР}, лаборатория микроэлектроники Института 

общей физики АН \mbox{СССР} и~Всесоюзный на\-уч\-но-ис\-сле\-до\-ва\-тель\-ский 

центр по изуче\-нию свойств по\-верх\-ности и~вакуума Госстандарта и~АН 

\mbox{СССР} (Москва), Институт проб\-лем технологии мик\-ро\-элект\-ро\-ни\-ки 

и~особо чис\-тых материалов с~СКБ (Черноголовка), Институт физики 

твер\-до\-го тела (Черноголовка), Институт проб\-лем информатики (15~тыс.~кв.\,м производственной площади и~5~тыс.~кв.\,м жилой), Саратовский 

филиал ИРЭ.

   

\section{40-летие Института проблем информатики}

   

   Постановление ЦК КПСС и~Совета Министров \mbox{СССР}, 

разрешающее АН СССР создать Институт проб\-лем информатики АН 

\mbox{СССР} (\mbox{ИПИАН}), было принято~29~июля 1983~г.~[6], 

а~уже~2~августа 1983~г.\ вы\-шло со\-от\-вет\-ст\-ву\-ющее распоряжение Президиума АН 

\mbox{СССР}. 

   

   В начале 1980-х~гг.\ появляются и~стремительно завоевывают новые 

области применения массовые средства вы\-чис\-ли\-тель\-ной техники~--- 

персональные компьютеры, существенное развитие получают научные 

осно\-вы информатики. Однако <<ве\-дом\-ст\-вен\-ность>> организаций, 

за\-ни\-ма\-ющих\-ся проб\-ле\-ма\-ми информатизации, стала тормозом для развития 

информатики в~\mbox{СССР}. Именно с~\mbox{целью} преодоления 

<<ведомственных барьеров>>, а~так\-же развития методов и~средств 

информатизации было образовано новое отделение АН \mbox{СССР},  

а~член-кор\-рес\-пон\-дент АН \mbox{СССР} Б.\,Н.~Наумов, возглавлявший 

Институт электронных управ\-ля\-ющих машин Минприбора \mbox{СССР}, 

был на\-зна\-чен ди\-рек\-то\-ром-ор\-га\-ни\-за\-то\-ром одного из новых институтов~--- 

\mbox{ИПИАН}. Основ\-ная задача ИПИАН была определена как <<проведение 

фундаментальных и~при\-клад\-ных исследований в~об\-ласти технических 

и~про\-грам\-мных средств массовой вы\-чис\-ли\-тель\-ной техники и~сис\-тем на их 

основе>>, а~интеллектуальной базой пер\-вых научных под\-раз\-де\-ле\-ний 

\mbox{ИПИАН} стали коллективы ряда научных отделов \mbox{ИНЭУМ}, 

переведенные в~\mbox{ИПИАН} в~начале~1984~г.

   

   Институт в~первые годы своего существования быстро развивался~--- в~его 

состав вошли филиалы в~Бердянске, Казани, Орле. В~1990~г.\ был образован 

совместный отдел с~радиотехническим институтом в~г.~Таганроге. Общая 

чис\-лен\-ность со\-труд\-ни\-ков института в~1988--1989~гг.\ превышала 

1000~человек. В~1989~г.\ директором института стал  

член-кор\-рес\-пон\-дент АН \mbox{СССР} Игорь Александрович Мизин 

(в~1997~г.\ из\-бран действительным членом РАН). В~1999~г.\ директором 

института стал Игорь Анатольевич Соколов (в~2003~г.\ избран  

чле\-ном-кор\-рес\-пон\-ден\-том РАН, а~в~2008~г.~--- действительным 

членом РАН). В~1992~г.\ институт получил новое наименование~--- 

Институт проб\-лем информатики Российской академии наук (ИПИ РАН). 

В~на\-сто\-ящее время преемником Института проб\-лем информатики РАН 

стал Федеральный исследовательский центр <<Информатика  

и~управ\-ле\-ние>> Российской академии наук (ФИЦ ИУ РАН)~--- 

подведомственная Министерству науки и~высшего образования Российской 

Федерации научная организация, вы\-пол\-ня\-ющая фундаментальные, 

поисковые и~при\-клад\-ные научные исследования и~разработки в~об\-ласти 

вычислительной и~при\-клад\-ной математики, сис\-тем\-но\-го анализа 

и~управ\-ле\-ния, тео\-ре\-ти\-че\-ской информатики и~информационных технологий, 

развития ин\-фор\-ма\-ци\-он\-но-те\-ле\-ком\-му\-ни\-ка\-ци\-он\-ной 

инфраструктуры и~информатизации общества. ФИЦ ИУ РАН был создан 

в~декабре 2014~г.\ путем реорганизации Федерального государственного 

бюд\-жет\-но\-го учреж\-де\-ния науки Института проб\-лем информатики 

Российской академии наук (ИПИ РАН) в~форме при\-со\-еди\-не\-ния к~нему 

Федерального государственного бюджетного учреж\-де\-ния науки Института 

сис\-тем\-но\-го анализа Российской академии наук (ИСА РАН) и~Федерального 

государственного бюджетного учреж\-де\-ния науки Вы\-чис\-ли\-тель\-но\-го цент\-ра 

им.\ А.\,А.~Дородницына Российской академии наук (ВЦ РАН). Центр имеет 

филиалы в~г.~Орле и~в~г.~Калининграде, которые до реорганизации были 

филиалами ИПИ РАН.  

На\-уч\-но-ме\-то\-ди\-че\-ское руководство ФИЦ ИУ РАН осуществляется 

двумя отделениями РАН~--- Отделением нанотехнологий 

и~информационных технологий (до 2002~г.~--- Отделение информатики, 

вы\-чис\-ли\-тель\-ной техники и~автоматизации, в~2002--2007~гг.~--- Отделение 

информационных технологий и~вы\-чис\-ли\-тель\-ных сис\-тем) и~Отделением 

математических наук. 

   

{\small\frenchspacing

{%\baselineskip=10.8pt

\addcontentsline{toc}{section}{Литература}

\begin{thebibliography}{9}

  \bibitem{1-zah}

  \Au{Zakharov~V.\,N.} Two fates in the history of computer technology in the USSR (Lebedev 

and Brouk)~/\!\!/ 4th Conference (International) on Computer Technology in Russia and in the 

Former Soviet Union Proceedings~/ Eds. I.~Kray\-ne\-va, A.~To\-mi\-lin.~--- Piscataway, NJ, USA: IEEE, 2018. 

P.~1--4. doi: 10.1109\!/\!SoRuCom.2017.00006.

  \bibitem{2-zah}

  Отечественная электронная вы\-чис\-ли\-тель\-ная техника~/ Сост. С.\,А.~Муравьев; под ред.\ С.\,В.~Хохлова.~--- М.: Столичная энциклопедия, 2014. 400~с.

  \bibitem{3-zah}

  \Au{Zakharov~V.} Computers and their application in the USSR in the middle of the 1980s: 

Situation, actions taken, predictions of development~/\!\!/ 3rd Conference (International) on 

Computer Technology in Russia and in the Former Soviet Union Proceedings.~--- Piscataway, NJ, USA: 

IEEE, 2014. P.~53--60. doi: 10.1109\!/\!SoRuCom.2014.19.

  \bibitem{4-zah}

  О дальнейшем развитии работ в~об\-ласти вы\-чис\-ли\-тель\-ной техники: По\-ста\-нов\-ле\-ние ЦК 

\mbox{КПСС} и~СМ \mbox{СССР} от 29.07.1983 №\,729-231.

  \bibitem{5-zah}

  О мерах по обеспечению работ в~об\-ласти вы\-чис\-ли\-тель\-ной техники и~ее применения 

в~народном хозяйстве: По\-ста\-нов\-ление СМ \mbox{СССР} от 29.07.1983 №\,730-232.

  \bibitem{6-zah}

  \Au{Захаров~В.\,Н.} Работы Института проб\-лем информатики РАН~/\!\!/ История 

отечественной электронной вы\-чис\-ли\-тель\-ной техники~/ Под ред. А.\,С.~Якунина.~--- М: 

Столичная энциклопедия, 2014. C.~444--448.

   \end{thebibliography}

} }







\vspace*{-6pt}



\hfill{\small\textit{Поступила в~редакцию 31.07.23}}



\vspace*{6pt}



\hrule



\vspace*{2pt}



\hrule



\vspace*{4pt}



%\newpage



%\vspace*{-28pt}







\def\tit{TO THE 40TH ANNIVERSARY OF~THE~DECREES OF~THE~CENTRAL 

COMMITTEE OF~THE~CPSU AND~THE~COUNCIL OF~MINISTERS OF~THE~USSR 

ON~THE~FURTHER DEVELOPMENT OF~WORK IN~THE~FIELD~OF~COMPUTER TECHNOLOGY, 

INCLUDING~AT~THE~USSR~ACADEMY~OF~SCIENCES}







\def\titkol{To the 40th Anniversary of~the~Decrees of~the~Central 

Committee of~the~CPSU}

% and~the~Council of~Ministers of~the~USSR 

%on~the~further development of~work in~the~field of~computer technology, 

%including at~the~USSR Academy of Sciences}



\def\aut{V.\,N.~Zakharov}



\def\autkol{V.\,N.~Zakharov}



\titel{\tit}{\aut}{\autkol}{\titkol}



\vspace*{-15pt}



\noindent

Federal Research Center ``Computer Science and Control'' of the Russian Academy of Sciences, 

44-2~Vavilov Str., Moscow 119133, Russian Federation





\def\leftfootline{\small{\textbf{\thepage}

\hfill Sistemy i Sredstva Informatiki~---

Systems and Means of Informatics 2023 vol~33 no\,3}

}%

 \def\rightfootline{\small{Sistemy i Sredstva Informatiki~---

 Systems and Means of Informatics 2023 vol~33 no\,3

\hfill \textbf{\thepage}}}



\vspace*{3pt} 





\Abste{The year 2023 marks the 40th anniversary of the establishment of the Institute of Problems of Informatics 

of the USSR Academy of Sciences, currently part of the Federal Research Center ``Computer Science and Control'' 

of the Russian Academy of Sciences. The Institute was established in accordance with a~decree of the Central 

Committee of the CPSU and the Council of Ministers of the USSR and a~decree of the Council of Ministers 

of the USSR dated July~29, 1983. The article discusses some prehistory of the development of computer technology 

in the USSR before these decrees were issued, the return of the subject of work on computer technology 

and computer science to the USSR Academy of Sciences, the esteblishment of the Department of Computer

Science, 

Computer Technology, and Automation of the USSR Academy of Sciences, and the main provisions of the decrees. 

Brief information about the development of IPIAN from the moment of its creation to the present time is also provided.}



\KWE{computer engineering; computer science; Academy of Sciences of the USSR; Central Committee of 

the CPSU; Council of Ministers of the USSR; Institute of Informatics Problems}



\DOI{10.14357\!/\!08696527230301}{XIWGAG}  



\vspace*{-12pt}





%\Ack

%\noindent



%\vspace*{-3pt}



\renewcommand{\bibname}{\protect\rmfamily References}

%\renewcommand{\bibname}{\large\protect\rm References}



{\small\frenchspacing

{ %\baselineskip=12pt

\addcontentsline{toc}{section}{References}

\begin{thebibliography}{9}

\bibitem{1-zah-1}

\Aue{Zakharov, V.\,N.} 2018. Two fates in the history of computer technology in the USSR (Lebedev and 

Brouk). \textit{4th Conference (International) on Computer Technology in Russia and in the Former Soviet 

Union Proceedings}. Eds.\ I.~Kray\-ne\-va and A.~To\-mi\-lin. Piscataway, NJ: IEEE. 1--4. doi: 

10.1109\!/\!SoRuCom.2017.00006.

\bibitem{2-zah-1}

Murav'ev, S.\,A., compiler. 2014.

\textit{Ote\-chest\-ven\-naya elekt\-ron\-naya vy\-chis\-li\-tel'\-naya tekh\-ni\-ka: Bio\-gra\-fi\-che\-skaya  

en\-tsik\-lo\-pe\-diya} [Domestic electronic computer technology: Biographical encyclopedia]. Ed.\ S.\,V.~Khokhlov. 

Moscow: Capital Encyclopedia. 400~p.

\bibitem{3-zah-1}

\Aue{Zakharov, V.\,N.} 2014. Computers and their application in the USSR in the middle of the 1980s: 

Situation, actions taken, predictions of development. \textit{3rd Conference (International) on Computer 

Technology in Russia and in the Former Soviet Union Proceedings}. Piscataway, NJ: IEEE. 53--60. doi: 

10.1109\!/\!SoRuCom.2014.19.

\bibitem{4-zah-1}

 O~dal'neyshem razvitii ra\-bot v~ob\-lasti vy\-chis\-li\-tel'\-noy tekh\-ni\-ki: Po\-sta\-nov\-le\-nie TsK 

KPSS i~SM SSSR ot 29.07.1983 No.\,729-231 [On the prospects for the development of work in the field of 

computer technology: Decree of the Central Committee of the CPSU and the Council of Ministers of the 

USSR No.\,729-231 dated 29.07.1983].

\bibitem{5-zah-1}

O~me\-rakh po obes\-pe\-che\-niyu ra\-bot v~ob\-lasti vy\-chis\-li\-tel'\-noy tekh\-ni\-ki i~ee  

pri\-me\-ne\-niya v~na\-rod\-nom kho\-zyay\-st\-ve: Po\-sta\-nov\-le\-nie SM SSSR ot 29.07.1983  

No.\,730-232 [On measures to ensure work in the field of computer technology and its application in the 

national economy: Decree of the Council of Ministers of the USSR No.\,730-232 dated 29.07.1983].

\bibitem{6-zah-1}

\Aue{Zakharov, V.\,N.} 2014. Ra\-bo\-ty In\-sti\-tu\-ta prob\-lem in\-for\-ma\-ti\-ki RAN [Proceedings of the 

Institute of Informatics Problems of the Russian Academy of Sciences]. \textit{Is\-to\-riya  

ote\-chest\-ven\-noy elekt\-ron\-noy vy\-chis\-li\-tel'\-noy tekh\-ni\-ki} [History of domestic computing 

technology]. Ed.\ A.\,S.~Yaku\-nin. Moscow: Capital Encyclopedia. 444--448. 

\end{thebibliography}

} }



\vspace*{-6pt}



\hfill{\small\textit{Received July 31, 2023}} 



\vspace*{-12pt} 





\Contrl



\vspace*{-3pt}



\noindent

\textbf{Zakharov Victor N.} (b.\ 1948)~--- Doctor of Science in technology, associate professor, scientific 

secretary, Federal Research Center ``Computer Science and Control'' of the Russian Academy of Sciences, 

44-2~Vavilov Str., Moscow 119333, Russian Federation; \mbox{vzakharov@ipiran.ru}

  



\label{end\stat}



\renewcommand{\bibname}{\protect\rm Литература}   

